\chapter{Crecimiento y \emph{moats}}\label{ch:growth}

% Alcance: la matriz de Ansoff, la taxonomía de moats (costos de cambio, efectos de red,
%   escala, intangibles), mecánica de los efectos de red, contabilidad del crecimiento.
% Vínculo cruzado: los efectos de red son el moat estructural; el volante operativo que los
%   genera está en \cref{ch:digital-funnel-and-crm}. Expansión de producto-mercado <-> STP
%   (\cref{ch:segmentation-targeting-positioning}).

Un diferencial de ventaja positivo (\cref{eq:advantage-spread}) responde si se tiene un
\emph{moat}; no dice cómo se creció hasta él ni cómo se protege y amplía. La estrategia de
\emph{growth} es el conjunto de decisiones sobre qué mercados y productos abordar a
continuación; la estrategia de moat es el conjunto de decisiones estructurales que hacen
defensible la posición resultante. Ambas operan a través del mismo cuello de botella
financiero: cada movimiento de \emph{growth} debe generar un NPV de cohorte que supere su
costo.

\section{La Matriz de Ansoff}\label{sec:ansoff-matrix}
% complexity: 3

¿Cuál de los cuatro vectores de crecimiento --- vender más a clientes existentes, lanzar
nuevos productos, entrar en nuevos mercados o diversificarse --- ofrece el mejor retorno por
unidad de riesgo de ejecución? El crecimiento siempre implica una apuesta sobre algo
desconocido. La matriz de Ansoff hace explícita esa apuesta: se apuesta por un producto
conocido en un mercado conocido (riesgo bajo), un nuevo producto en un mercado conocido,
un producto conocido en un nuevo mercado, o un nuevo producto en un nuevo mercado (riesgo
alto). El riesgo aumenta en diagonal, y también el retorno requerido.

\textcite{kotler2021} presentan la matriz como una estructura de decisión \(2\times2\)
según la familiaridad con el producto (existente / nuevo) y la familiaridad con el mercado
(existente / nuevo):
\begin{description}
  \item[Penetración de mercado] Producto existente, mercado existente. Aumentar cuota
        mediante precio, promoción o distribución. Menor riesgo; mayor eficiencia de
        recursos; techo limitado en un mercado maduro.
  \item[Desarrollo de producto] Nuevo producto, mercado existente. Extender la plataforma
        para clientes que la empresa ya conoce. Riesgo moderado (incertidumbre de producto,
        pero se conserva el conocimiento del cliente).
  \item[Desarrollo de mercado] Producto existente, nuevo mercado. Replicar la propuesta de
        valor en una nueva geografía, vertical o segmento. Riesgo moderado (producto
        conocido, cliente desconocido). Requiere una nueva decisión de STP
        (\cref{ch:segmentation-targeting-positioning}).
  \item[Diversificación] Nuevo producto, nuevo mercado. El cuadrante de mayor riesgo;
        justificado sólo cuando los mercados principales están saturados o en declive, y
        generalmente ejecutado vía M\&A en lugar de construcción orgánica.
\end{description}

«Mercado» y «producto» son elecciones definicionales, no hechos objetivos. El mismo
movimiento puede etiquetarse como penetración («sólo estamos añadiendo funciones») o
desarrollo de producto («es una nueva línea de producto») según el encuadre. La versión útil
de la matriz obliga a un juicio concreto: ¿requiere la nueva oferta capacidades o
movimientos de GTM que la empresa no tiene actualmente? Si es así, es un movimiento de
nuevo producto o nuevo mercado; si no, es penetración. Sin ese rigor la matriz se convierte
en vocabulario, no en herramienta.

\begin{example}
\emph{Decisión: ¿qué vector de crecimiento debe priorizar la empresa SaaS en los próximos
18 meses?}

Estado actual: \money{4000000} ARR, \num{200} clientes nuevos/mes, \money{340000} MRR.
Las cuatro opciones evaluadas contra el NPV de cohorte:

\begin{itemize}
  \item \textbf{Penetración de mercado}: el \emph{funnel} convierte al \qty{2.5}{\percent}
        (200/8000 sesiones calificadas). Hacer pruebas A/B en \emph{landing pages} y
        mejorar el paso de conversión de prueba a pago (\cref{ch:digital-funnel-and-crm})
        podría elevar la conversión al \qty{3}{\percent} --- 240 pagos/mes --- a un costo
        incremental casi nulo. El NPV más alto por dólar invertido a la escala actual.
  \item \textbf{Desarrollo de producto}: un complemento de informes a \money{20}/mes
        elevaría el ARPU de \money{50} a \money{70}, aumentando el LTV de \money{3150} a
        \(\money{70}\times 0.42/(0.11-0.03) = \money{4410}\). La restricción es si los
        clientes existentes lo desean
        (\cref{ch:segmentation-targeting-positioning}).
  \item \textbf{Desarrollo de mercado}: entrar en el segmento europeo de pymes requiere
        cumplimiento con GDPR, métodos de pago locales y un flujo de \emph{onboarding}
        localizado --- costo de ejecución significativo con CAC incierto.
  \item \textbf{Diversificación}: prematuro con \money{4000000} ARR; las reservas de
        atención ejecutiva son la restricción vinculante.
\end{itemize}

La penetración y el desarrollo de producto son los primeros movimientos racionales; el
desarrollo de mercado es una opción a 24 meses una vez demostradas las economías unitarias.
\end{example}

\begin{admonition}{Advertencia}
Usar la matriz de Ansoff como menú en lugar de herramienta de secuenciación. Las cuatro
opciones siempre están disponibles en teoría; el valor de la matriz radica en forzar un
ranking explícito por retorno ajustado al riesgo y capacidad de ejecución. Perseguir dos
cuadrantes de alto riesgo simultáneamente diluye la atención gerencial por debajo de la
dosis mínima eficaz en cualquiera de ellos.
\end{admonition}

La matriz de Ansoff es el análogo corporativo-estratégico de la decisión de segmentación en
\cref{ch:segmentation-targeting-positioning}: ambas plantean la misma pregunta (¿a quién
servimos y con qué?), una a nivel de portafolio y la otra a nivel de producto. La
investigación en estrategia corporativa (\parencite{dranove7theconomics}) concluye que la
diversificación \emph{relacionada} --- moverse a lo largo de una sola dimensión de la matriz
a la vez --- genera retornos superiores a la diversificación no relacionada, consistente con
el argumento VRIN de que la ventaja es específica al recurso y no se transfiere entre
negocios no relacionados.

La secuencia de Ansoff debe seguir la taxonomía de moats más adelante: penetración primero,
desarrollo de producto una vez establecido el moat, desarrollo de mercado una vez probado el
producto. \textcite{kotler2021} contextualizan la matriz dentro de la estrategia moderna de
\emph{growth}.

\section{La Taxonomía de Moats}\label{sec:moat-taxonomy}
% complexity: 3

¿Qué barrera estructural protegerá los retornos de la empresa una vez que un competidor
tenga los recursos para igualar la calidad del producto y el precio? El moat debe construirse
antes de que sea necesario. Un moat es cualquier ventaja estructural que hace irracional para
un competidor bien capitalizado atacar directamente. A diferencia de una ventaja de producto,
un moat es auto-reforzante: cuantos más clientes se sirven, más difícil resulta para un
rival desplazarlos. Los moats convierten una ventaja temporal en un diferencial de valor
positivo sostenido (\cref{eq:advantage-spread}).

Cuatro tipos de moat sostienen ROIC $>$ WACC (\cref{eq:roic}) en un horizonte estratégico
\parencite{porter1985,dranove7theconomics}:
\begin{description}
  \item[Costo de cambio] El costo directo e indirecto que un cliente incurre al migrar a
        un rival. Incluye migración de datos, reentrenamiento, interrupción de flujos de
        trabajo y reconstrucción de integraciones. El costo de cambio eleva el precio
        efectivo de salida sin elevar el precio nominal de permanencia --- amplía la banda
        de precios dentro de la cual el titular está protegido de la deserción.
  \item[Efecto de red] Cada usuario adicional hace el producto más valioso para todos los
        demás usuarios. Este es el moat más poderoso en los mercados tecnológicos
        (formalizado en \cref{sec:network-effects}): es escala del lado de la demanda, no
        del lado del costo, y se compone automáticamente.
  \item[Escala del lado del costo] Reducción de costos unitarios a gran escala --- dilución
        del costo fijo, apalancamiento de compras o efectos de la curva de aprendizaje.
        Protege vía precio: los titulares pueden subcotizar rentablemente a un entrante que
        aún no ha alcanzado la escala mínima eficiente.
  \item[Activos intangibles] Patentes, licencias regulatorias o reputación de marca que los
        competidores no pueden replicar rápidamente. La marca es el moat intangible más
        débil en B2B SaaS (los compradores son profesionales y sensibles al precio); las
        patentes requieren una inversión densa en propiedad intelectual.
\end{description}

Los moats se erosionan. Los costos de cambio colapsan cuando emerge una capa estándar de
portabilidad de datos (por ejemplo, el derecho a la portabilidad de datos del GDPR en SaaS).
Los efectos de red se revierten cuando una plataforma alcanza deseconomías de escala (spam,
congestión). La escala del lado del costo es socavada por una tecnología disruptiva que
resetea la estructura de costos. Tratar un moat como un activo en depreciación que requiere
reinversión, no como una barrera permanente.

\begin{example}
\emph{Decisión: ¿cuál es el moat principal de la empresa SaaS y cuánto vale?}

Moat principal: \emph{costo de cambio} derivado de la profundidad de las integraciones. Un
cliente con 10 integraciones activas enfrenta un costo de migración estimado de
\money{3000}--\money{6000} (tiempo de ingeniería, reentrenamiento, riesgo de tiempo de
inactividad) --- muy por encima del valor anual de la suscripción de \money{600}. Esto
significa que el costo de cambio efectivo del cliente supera un año completo de ingresos, lo
que suprime materialmente el \emph{churn} por debajo del promedio estructural.

Expresión financiera: si el moat de costo de cambio reduce el \emph{churn} anual del
\qty{10}{\percent} al \qty{7}{\percent}, la tasa mensual neta de \emph{churn} se desplaza
del \qty{0.87}{\percent} al \qty{0.60}{\percent} y el LTV aumenta:
\[
  \Delta\mathrm{LTV} \approx \frac{\money{252}}{0.08} - \frac{\money{252}}{0.11}
                           = \money{3150} - \money{2291} \approx \money{860}.
\]
Un incremento de \money{860} en LTV significa que la empresa puede gastar \money{860} más
por cliente adquirido sin violar el piso LTV:CAC (\cref{eq:ltv-cac}), lo que a 200
clientes/mes desbloquea \money{172000}/mes de presupuesto incremental de CAC ---
\qty{143}{\percent} del gasto publicitario actual.
\end{example}

\begin{admonition}{Advertencia}
Llamar \emph{retention} a un moat. La \emph{retention} es el síntoma; el moat es la causa
estructural. Un producto que retiene clientes porque es excelente los pierde el día que un
competidor construye un producto mejor. Un producto que retiene clientes porque las
integraciones hacen dolorosa la migración los retiene incluso cuando existe un producto
mejor. Invertir en la causa, no en el síntoma.
\end{admonition}

El análisis de moats conecta con la literatura de teoría de precios a través del índice de
Lerner (\cref{eq:lerner}): una empresa con costos de cambio puede fijar precio por encima
del costo marginal en una cantidad proporcional al costo de cambio dividido por el horizonte
del contrato. El moat es, en términos de teoría de precios, una fuente de inelasticidad:
reduce la elasticidad precio propia de la demanda (\cref{eq:elasticity}) del producto del
titular sin cambiar la elasticidad de la categoría en su conjunto.

Los costos de cambio son el moat principal de la empresa SaaS; los efectos de red
(\cref{sec:network-effects}) son el moat secundario potencial si se construye un mercado de
integraciones. Juntos generan el diferencial de valor sostenido que define
\cref{ch:advantage-and-disruption}. \textcite{porter1985} y \textcite{dranove7theconomics}
son las fuentes principales.
\textcite{helmer2016powers} cataloga los mismos mecanismos como siete \emph{powers}: economías de escala, economías de red, contraposicionamiento, costos de cambio, marca, recurso acaparado y \emph{process power} --- cada uno una propiedad estructural que sostiene el diferencial ROIC--WACC.

\section{Mecánica de los Efectos de Red}\label{sec:network-effects}
% complexity: 4 -> reutilizar figura growth-loop

¿El producto tiene efectos de red, qué tan fuertes son y qué implica eso para el
presupuesto de CAC y la trayectoria de \emph{growth}? Los efectos de red crean un moat del
lado de la demanda que se compone: cada nuevo usuario eleva el valor del producto para todos
los usuarios existentes, lo que facilita la venta al siguiente usuario y reduce el CAC con
el tiempo. La implicación estructural es que las economías unitarias \emph{mejoran} con la
escala --- lo opuesto a la mayoría de las estructuras de costo.

La ley de Metcalfe aproxima el valor de una red de \(n\) nodos como proporcional al número
de pares distintos:
\begin{equation}\label{eq:metcalfe}
  V(n) \;\propto\; n(n-1) \;\approx\; n^2 \quad\text{for large }n.
\end{equation}
La escala cuadrática es el núcleo económico: doblar usuarios más que dobla el valor, por lo
que el valor por usuario aumenta con el tamaño de la red. Esto significa que la disposición
a pagar aumenta a medida que la red crece, sin ningún cambio en el producto mismo --- la
curva de demanda se desplaza hacia afuera únicamente por el tamaño de la base instalada.

\begin{figure}[htbp]
  \centering
  \includegraphics[width=0.86\linewidth]{growth-loop}
  \caption{El bucle de crecimiento auto-reforzante. Los nuevos clientes producen referidos,
    contenido o datos que atraen a más clientes nuevos. Cuando el coeficiente viral
    \(k > 1\) (\cref{eq:viral-k}), el bucle es autosustentable y el CAC de los canales
    pagados cae a medida que aumenta la adquisición orgánica.}
  \label{fig:network-effects}
\end{figure}

La \cref{eq:metcalfe} supone que todos los vínculos tienen igual valor, lo que es
empíricamente falso. Las redes reales siguen una distribución de ley de potencia: un número
reducido de nodos de alto valor (usuarios intensivos, cuentas empresariales) contribuyen de
forma desproporcionada al valor de la red. La fórmula cuadrática sobreestima el valor en
redes escasas y heterogéneas. También confunde los efectos de red \emph{directos} (el
usuario se beneficia de otros usuarios directamente) con los \emph{indirectos} (el usuario
se beneficia de los complementos atraídos por una base de usuarios grande). Los efectos
directos son más fuertes pero más raros; los indirectos son más comunes en los mercados de
plataformas (\cref{ch:business-models}).

\begin{example}
\emph{Decisión: si la empresa SaaS construye un mercado de integraciones que conecta
proveedores, ¿cuál es el valor potencial con \(n = 100\) integraciones?}

A partir de \cref{eq:metcalfe}, el número de pares distintos de proveedores con \(n = 100\):
\[
  \frac{n(n-1)}{2} = \frac{100 \times 99}{2} = 4{,}950 \text{ pairs}.
\]
Cada par representa una oportunidad potencial de intercambio de datos o automatización de
flujos de trabajo que no existía con \(n = 50\) (que sólo ofrecía \num{1225} pares). El
valor \emph{marginal} de la integración número 100 son \(99\) pares nuevos --- 99 veces el
valor de la primera. Por eso los mercados de integraciones justifican una proporción
desproporcionada de la inversión en ingeniería en las etapas tempranas: el valor de opción
del nodo número 100 está embebido en la decisión del primer nodo de unirse.

El bucle de crecimiento (\cref{fig:network-effects}) captura la expresión organizacional:
cada integración atrae a más proveedores, cuya participación atrae a más clientes, cuyos
datos mejoran el producto, que atrae a más integraciones. El coeficiente viral
(\cref{eq:viral-k}) es la instanciación del lado de la demanda del mismo bucle aplicado a
usuarios en lugar de integraciones. Ambos son canales hacia \cref{ch:digital-funnel-and-crm}.
\end{example}

\begin{admonition}{Advertencia}
Tratar efectos débiles del mismo lado como efectos de red a escala Metcalfe. Un producto
SaaS utilizado de forma independiente por cada cliente (sin datos compartidos, sin
interacciones entre usuarios) tiene efecto de red directo igual a cero independientemente
del número de usuarios. La prueba relevante es: ¿cambia la experiencia del usuario \(j\)
cuando se une el usuario \(k\)? Si no es así, no es un efecto de red; es una característica
del producto que simplemente escala. Sobrestimar los efectos de red distorsiona la
asignación de capital.
\end{admonition}

Los efectos de red directos frente a los indirectos tienen implicaciones diferentes para
la masa crítica. Los efectos directos requieren una red mínima viable de usuarios activos
antes de que el valor supere al del producto independiente; los efectos indirectos requieren
un conjunto mínimo viable de complementadores. El problema del huevo y la gallina en el
lanzamiento de plataformas (\cref{ch:business-models}) es la versión del lado de la demanda
del requisito de masa crítica en \cref{eq:metcalfe}.
\parencite{dranove7theconomics} desarrollan las implicaciones de bienestar y precios.

Los efectos de red son el moat más sólido en los mercados tecnológicos pero el más difícil
de construir desde cero. El coeficiente viral (\cref{eq:viral-k}) en
\cref{ch:digital-funnel-and-crm} es el mecanismo del lado de la demanda; las métricas para
su seguimiento están en \cref{ch:metrics-catalog}. \textcite{dranove7theconomics} aportan
el fundamento teórico.

\section{Contabilidad del Crecimiento}\label{sec:growth-accounting}
% complexity: 4 -> figura mrr-waterfall

¿La empresa realmente está creciendo, o los ingresos nuevos simplemente reemplazan los
ingresos perdidos por \emph{churn}? La descomposición del crecimiento neto de MRR indica en
qué canales invertir y dónde están las fugas. El crecimiento agregado de ingresos confunde
cuatro actividades empresariales muy distintas. La adquisición de nuevos clientes, la
expansión de cuentas existentes, la contracción de cuentas que bajan de plan y el
\emph{churn} de cuentas perdidas se suman a un único titular que enmascara las dinámicas
subyacentes. Una empresa con ventas nuevas sólidas y \emph{churn} alto parece idéntica en
el titular a una empresa con ventas nuevas modestas y \emph{retention} sólida, pero
requieren respuestas estratégicas completamente distintas.

Sea el MRR al inicio de un período \(M_0\). El MRR al cierre es:
\begin{equation}\label{eq:mrr-growth}
  M_1 \;=\; M_0 \;+\; \Delta_{\text{new}}
               \;+\; \Delta_{\text{expansion}}
               \;-\; \Delta_{\text{contraction}}
               \;-\; \Delta_{\text{churn}},
\end{equation}
donde cada \(\Delta\) es un número positivo. El nuevo MRR neto es
\(\Delta_{\text{new}} + \Delta_{\text{expansion}} - \Delta_{\text{contraction}} - \Delta_{\text{churn}}\).
La Retención Neta de Ingresos (NRR), derivada en \cref{ch:business-models}, es:
\[
  \mathrm{NRR} = \frac{M_0 + \Delta_{\text{expansion}} - \Delta_{\text{contraction}} - \Delta_{\text{churn}}}{M_0},
\]
es decir, cómo se ve la \emph{retention} de ingresos \emph{excluyendo} a los nuevos
clientes. NRR $>$ 100\% significa que la base existente se expande en términos netos ---
la empresa podría teóricamente cesar las ventas nuevas y seguir creciendo.

\begin{figure}[htbp]
  \centering
  \includegraphics[width=0.86\linewidth]{mrr-waterfall}
  \caption{Cascada de MRR para el ejemplo SaaS de referencia. El MRR inicial de
    \money{340000} crece por clientes nuevos (\money{30000}) y expansión (\money{20000}),
    y se reduce por contracción (\money{5000}) y \emph{churn} (\money{22000}), terminando
    en \money{363000} --- una ganancia neta del \qty{6.8}{\percent}.}
  \label{fig:mrr-waterfall}
\end{figure}

La \cref{eq:mrr-growth} es una identidad --- siempre se cumple por construcción. El poder
diagnóstico proviene de evaluar si cada componente es saludable \emph{en relación con el
modelo de negocio}. Un producto transaccional de alta velocidad (SaaS de comercio
electrónico) se espera que tenga mayor \emph{churn} bruto y menor expansión; un producto
empresarial de expansión de asientos se espera que tenga alta expansión y bajo \emph{churn}
bruto. Comparar niveles absolutos de \emph{churn} sin ajustar por tipo de modelo de negocio
produce referencias falsas. En relación con esto, la ecuación trata todo el MRR como
igualmente valioso; un dólar de \emph{churn} de un cliente con 36 meses de antigüedad es
más preocupante que uno de una cuenta de un mes (\cref{eq:clv}).

\begin{example}
\emph{Decisión: ¿dónde debe invertir el equipo --- en adquisición nueva o en reducción del
\emph{churn}?}

De la hoja de datos: MRR = \money{340000}. Movimientos mensuales observados:
Movimientos mensuales observados (supuestos declarados; la combinación de planes y el empuje
de expansión elevan \(\Delta_{\text{new}}\) por encima del simple producto
\num{200}\(\times\)\money{50}):
\[
  \begin{aligned}
    \Delta_{\text{new}} &= +\money{30000} \quad (\text{new customers across plan mix})\\
    \Delta_{\text{expansion}} &= +\money{20000} \quad (\text{seat adds, plan upgrades})\\
    \Delta_{\text{contraction}} &= -\money{5000} \quad (\text{downgrades})\\
    \Delta_{\text{churn}} &= -\money{22000} \quad (\approx\qty{6.5}{\percent}\text{ of }\money{340000})
  \end{aligned}
\]
MRR neto nuevo: $\money{30000}+\money{20000}-\money{5000}-\money{22000}=+\money{23000}$
(+\qty{6.8}{\percent}/mes).
\[
  \mathrm{NRR} = \frac{\money{340000}+\money{20000}-\money{5000}-\money{22000}}{\money{340000}}
               = \frac{\money{333000}}{\money{340000}} \approx \qty{97.9}{\percent}.
\]
Un NRR del \qty{97.9}{\percent} indica que la base existente se contrae levemente sin ventas
nuevas --- la empresa es neta-negativa en \emph{retention}. Incluso mejoras modestas en la
\emph{retention} se componen dramáticamente: reducir el \emph{churn} de \money{22000} a
\money{15000}/mes eleva el NRR al \qty{99.7}{\percent} y aumenta el MRR al cierre en
\money{7000}/mes sin adquirir un solo cliente nuevo. Ese \money{7000} se compone. La
prioridad de inversión es la reducción del \emph{churn} --- específicamente los recursos
VRIN (integraciones) que elevan los costos de cambio (\cref{sec:moat-taxonomy}).
\end{example}

\begin{admonition}{Advertencia}
Optimizar el nuevo MRR \emph{bruto} ignorando el \emph{churn}. Una empresa que agrega
\money{40000}/mes en nuevo MRR mientras pierde \money{40000}/mes por \emph{churn} está
corriendo en el lugar: NRR $=$ \qty{100}{\percent} exactamente, pero la base de clientes se
ha renovado por completo en un horizonte finito y el CAC se ha gastado con retorno compuesto
nulo. El crecimiento \emph{neto} de MRR y el NRR son las métricas correctas; el nuevo MRR
bruto es un número de vanidad sin el contexto de la \emph{retention}.
\end{admonition}

La cascada de MRR es la versión operativa del modelo de LTV a nivel de cohorte en
\cref{ch:customer-economics}: la \cref{eq:clv} descuenta los flujos de caja esperados de
una sola cohorte; la \cref{eq:mrr-growth} es la manifestación mensual agregada de todas las
cohortes superpuestas sumadas en el tiempo calendario. El catálogo de métricas
(\cref{ch:metrics-catalog}) proporciona referencias para cada componente.
\textcite{ries2011} introdujo el encuadre de «contabilidad del crecimiento» para las
métricas de producto; \textcite{bendle2020marketing} lo extienden al contexto de
atribución de \emph{marketing}.

La contabilidad del crecimiento cierra el ciclo entre las decisiones de estrategia de
\emph{growth} (Ansoff) y las inversiones en moats (costos de cambio, efectos de red) al
hacer visible en tiempo real el resultado financiero de cada una. La \cref{eq:mrr-growth} y
el NRR son el panel de control principal para \cref{ch:metrics-catalog}.
\textcite{ries2011} y \textcite{bendle2020marketing} son las fuentes principales.
